% Unit 111 terjemahan Bahasa Indonesia dari chapter8.tex baris 925--1058.
% Sumber: Methods of Algebra, Volume 2, commit
% 9a5803ff77dd3257484cb177f851a73770a59dd3 (CC BY 4.0).
% stable-unit-id: o014.aljabr2.chapter8.dold-kan-b
% Cakupan: sisa Bagian 8.5, dari Teorema Dold--Kan hingga catatan mengenai
% versi kokompleks dan kosimplisial; Unit 112 memulai Bagian 8.6 pada baris 1059.

\hypertarget{o014.aljabr2.chapter8.dold-kan-b}{ }
% segment-id: o014.aljabr2.chapter8.dold-kan-b.t001
\begin{theorem}[A.\ Dold, D.\ Kan]\label{prop:Dold-Kan}
	\index{korespondensi Dold-Kan@korespondensi Dold--Kan (Dold--Kan correspondence)}
	Untuk setiap kategori aditif $\mathcal{A}$, funktor
	$\Gamma:\cate{Ch}_{\geq0}(\mathcal{A})\to\cate{s}\mathcal{A}$ merupakan
	funktor penuh dan setia.
	\begin{itemize}
		% segment-id: o014.aljabr2.chapter8.dold-kan-b.i001
		\item Jika $\mathcal{A}$ juga merupakan kategori Karoubi sebagaimana
		disebutkan dalam \S\ref{sec:direct-sum}, yaitu setiap idempoten mempunyai
		kernel, maka $\Gamma$ merupakan ekuivalensi kategori.
		% segment-id: o014.aljabr2.chapter8.dold-kan-b.i002
		\item Jika lebih lanjut $\mathcal{A}$ merupakan kategori abelian (dan
		karenanya juga kategori Karoubi), maka $\mathrm{N}$ adalah funktor
		kuasi-invers bagi $\Gamma$ dan
		% segment-id: o014.aljabr2.chapter8.dold-kan-b.g001
		$\begin{tikzcd}
			\Gamma: \cate{Ch}_{\geq 0}(\mathcal{A}) \arrow[shift left, r] & \cate{s}\mathcal{A} : \mathrm{N} \arrow[shift left, l]
		\end{tikzcd}$
		merupakan adjungsi yang sekaligus merupakan ekuivalensi. Morfisme unit dan kounit adjungsi
		tersebut dideskripsikan dalam Lema~\ref{prop:GammaN-adjunction}.
	\end{itemize}
\end{theorem}
% segment-id: o014.aljabr2.chapter8.dold-kan-b.pf001
\begin{proof}
	Tinjau funktor
	% segment-id: o014.aljabr2.chapter8.dold-kan-b.q001
	\[ \tilde{h}_{\mathcal{A}}: \mathcal{A} \to \tilde{\mathcal{A}}^\wedge := \cate{Ab}^{(\mathcal{A}^{\opp})}, \quad Y \mapsto \Hom_{\mathcal{A}}(\cdot, Y). \]
	Komposisinya dengan funktor pelupa $\cate{Ab}\to\cate{Set}$ sama dengan
	pembenaman Yoneda
	$h_{\mathcal{A}}:\mathcal{A}\to\mathcal{A}^\wedge$ yang dibahas
	dalam \S\ref{sec:Yoneda}. Funktor $\tilde{h}_{\mathcal{A}}$ penuh dan setia.
	Pernyataan ini tidak lain adalah Lema Yoneda versi diperkaya atas $\cate{Ab}$,
	tetapi juga dapat diturunkan dari versi semula. Untuk setiap
	$Y_1,Y_2\in\Obj(\mathcal{A})$, komposisi pemetaan
	% segment-id: o014.aljabr2.chapter8.dold-kan-b.q002
	\[ \Hom_{\mathcal{A}}(Y_1, Y_2) \to \Hom_{\tilde{\mathcal{A}}^\wedge}\left(\tilde{h}_{\mathcal{A}}(Y_1), \tilde{h}_{\mathcal{A}}(Y_2) \right) \to \Hom_{\mathcal{A}^\wedge}\left( h_{\mathcal{A}}(Y_1), h_{\mathcal{A}}(Y_2)\right) \]
	diketahui bijektif, sedangkan pemetaan kedua jelas injektif. Karena itu,
	kedua pemetaan tersebut bijektif.

	Proposisi~\ref{prop:functor-cat-Abel} menunjukkan bahwa
	$\tilde{\mathcal{A}}^\wedge$ secara alami merupakan kategori abelian. Gunakan
	lambang $\tilde{h}_{\mathcal{A}}$ yang sama bagi funktor yang diinduksinya
	pada kompleks rantai dan objek simplisial, lalu tinjau diagram
	% segment-id: o014.aljabr2.chapter8.dold-kan-b.e001
	\begin{equation*}\begin{tikzcd}
		\cate{Ch}_{\geq 0}(\mathcal{A}) \arrow[d, "{\tilde{h}_{\mathcal{A}}}"'] \arrow[r, "\Gamma"] & \cate{s}\mathcal{A} \arrow[d, "{\tilde{h}_{\mathcal{A}}}"] \\
		\cate{Ch}_{\geq 0}(\tilde{\mathcal{A}}^\wedge) \arrow[r, "\Gamma"'] & \cate{s}\tilde{\mathcal{A}}^\wedge.
	\end{tikzcd}\end{equation*}
	Diagram ini komutatif hingga isomorfisme kanonik. Pengamatan pada paragraf
	sebelumnya menunjukkan bahwa kedua anak panah vertikal penuh dan setia,
	sedangkan penerapan Lema~\ref{prop:DK-Ab} pada setiap objek menunjukkan
	bahwa baris kedua merupakan ekuivalensi kategori. Jadi, $\Gamma$ pada baris
	pertama penuh dan setia.

	Untuk setiap funktor $F:\mathcal{C}\to\mathcal{C}'$, suatu objek
	$X'\in\Obj(\mathcal{C}')$ dikatakan berada dalam \emph{citra esensial} $F$
	jika terdapat $X\in\Obj(\mathcal{C})$ dengan $X'\simeq FX$. Dipadukan
	dengan deskripsi kuasi-invers bagi
	$\cate{Ch}_{\geq0}(\cate{Ab})\to\cate{s}\cate{Ab}$ dalam
	Lema~\ref{prop:DK-Ab}, hal ini memberikan kesimpulan berikut:
	$X\in\Obj(\cate{s}\mathcal{A})$ berada dalam citra esensial $\Gamma$ jika
	dan hanya jika $(\mathrm{N}\tilde{h}_{\mathcal{A}}(X))_n$ berada dalam citra
	esensial $\tilde{h}_{\mathcal{A}}:\mathcal{A}\to
	\tilde{\mathcal{A}}^\wedge$ untuk setiap $n\geq0$.

	Perhatikan bahwa
	$\tilde{h}_{\mathcal{A}}:\mathcal{A}\to\tilde{\mathcal{A}}^\wedge$
	mempertahankan jumlah langsung hingga. Jika $\mathcal{A}$ kategori Karoubi,
	citra esensial $\mathcal{A}\to\tilde{\mathcal{A}}^\wedge$ karenanya
	tertutup terhadap pengambilan suku langsung. Karena
	$(\mathrm{N}\tilde{h}_{\mathcal{A}}(X))_n$ merupakan suku langsung dari
	$(\Gamma\mathrm{N}\tilde{h}_{\mathcal{A}}(X))_n\simeq
	\tilde{h}_{\mathcal{A}}(X)_n$,\footnote{Catatan penerjemah: ruas terakhir
	tercetak sebagai $h_{\mathcal{A}}(X)_n$ dalam sumber. Objek-objek lain dalam
	isomorfisme ini merupakan presheaf bernilai-$\cate{Ab}$, sehingga ruas yang
	bertipe benar adalah $\tilde{h}_{\mathcal{A}}(X)_n$.} dalam keadaan ini
	$\Gamma$ penuh, setia, dan surjektif esensial, sehingga merupakan
	ekuivalensi.

	Terakhir, teorema ekuivalensi adjoin \cite[Teorema 2.6.12]{Li1} memastikan
	bahwa suatu kuasi-invers bagi
	$\Gamma:\cate{Ch}_{\geq0}(\mathcal{A})\to\cate{s}\mathcal{A}$ selalu dapat
	dijadikan adjoin kanannya. Jika $\mathcal{A}$ kategori abelian,
	Lema~\ref{prop:GammaN-adjunction} telah menunjukkan bahwa $\mathrm{N}$
	adalah adjoin kanan $\Gamma$, sehingga juga merupakan kuasi-inversnya.
	Selesailah bukti.
\end{proof}

% segment-id: o014.aljabr2.chapter8.dold-kan-b.cv001
\begin{convention}
	Berdasarkan hasil ini, jika $\mathcal{A}$ kategori Karoubi, kita dapat
	memilih suatu funktor kuasi-invers bagi
	$\Gamma:\cate{Ch}_{\geq0}(\mathcal{A})\to\cate{s}\mathcal{A}$ dan secara
	wajar menuliskannya sebagai $\mathrm{N}$.
\end{convention}

% segment-id: o014.aljabr2.chapter8.dold-kan-b.p001
Walaupun $\mathrm{N}X$ mula-mula didefinisikan sebagai subobjek
$\mathrm{C}X$, ia juga dapat dipahami sebagai objek hasil bagi; sudut pandang
terakhir sering lebih mudah digunakan. Berikut penjelasannya.

% segment-id: o014.aljabr2.chapter8.dold-kan-b.p002
Tinjau kategori Karoubi $\mathcal{A}$ dan
$X\in\Obj(\cate{s}\mathcal{A})$. Untuk setiap $n\geq0$, definisikan morfisme
$\Psi_n$ dalam $\mathcal{A}$ sebagai morfisme yang dicirikan oleh
komutativitas diagram berikut untuk setiap $0\leq j<n$:
% segment-id: o014.aljabr2.chapter8.dold-kan-b.g002
\[\begin{tikzcd}
	\bigoplus_{0 \leq i < n} X_{n-1} \arrow[r, "\Psi_n"] & X_n \\
	X_j \arrow[hookrightarrow, u, "\text{suku langsung}"] \arrow[ru, "{s_j}"'] &
\end{tikzcd}\]
Jadi, mengambil $\Coker(\Psi_n)$ (yakni kopenyama dari $\Psi_n$ dan $0$)
secara intuitif menghapus semua bagian degenerat dari $X_n$, setidaknya
jika $\mathcal{A}$ kategori abelian. Hasil berikut bukan hanya menjamin bahwa
$\Coker(\Psi_n)$ ada, tetapi juga mengidentifikasikannya dengan
$(\mathrm{N}X)_n$ melalui penanaman alami
$u:\mathrm{N}X\to\mathrm{C}X$ dalam Proposisi~\ref{prop:Dold-Kan-u}. Ini
memberikan sudut pandang lain atas $\mathrm{N}X$.

% segment-id: o014.aljabr2.chapter8.dold-kan-b.pr001
\begin{proposition}\label{prop:Dold-Kan-v}
	Dalam keadaan di atas, $\Coker(\Psi_n)$ ada dan isomorfik secara kanonik
	dengan $(\mathrm{N}X)_n$. Tinjau komposisi yang bersesuaian
	$v_n:X_n\twoheadrightarrow\Coker(\Psi_n)\rightiso(\mathrm{N}X)_n$.
	Maka $(v_n)_n$ memberikan morfisme
	$v:\mathrm{C}X\to\mathrm{N}X$ yang memenuhi
	$vu=\identity_{\mathrm{N}X}$ untuk penanaman alami
	$u:\mathrm{N}X\hookrightarrow\mathrm{C}X$; konstruksi ini funktorial dalam $X$.
\end{proposition}
% segment-id: o014.aljabr2.chapter8.dold-kan-b.pf002
\begin{proof}
	Berdasarkan Teorema~\ref{prop:Dold-Kan}, kita dapat mengandaikan bahwa
	$X=\Gamma A$, dengan
	$A\in\Obj(\cate{Ch}_{\geq0}(\mathcal{A}))$. Dalam hal ini, $\Psi_n$
	dapat ditulis sebagai
	% segment-id: o014.aljabr2.chapter8.dold-kan-b.q003
	\[ \bigoplus_{0 \leq i < n} \bigoplus_{t: [n-1] \twoheadrightarrow [k]} A_k \to \bigoplus_{t': [n] \twoheadrightarrow [k] } A_k. \]
	Dari Definisi~\ref{def:DK-Gamma} tentang funktor $\Gamma$, terlihat bahwa
	pada suku langsung $(i,t)$ di ruas kiri, $\Psi_n$ bekerja sebagai berikut.
	Ambil komposisi
	$[n]\xrightarrow{\mathrm{s}^i}[n-1]\xrightarrow{t}[k]$, yang masih
	merupakan pemetaan monoton surjektif, dan tuliskan $t'$ untuk komposisi itu;
	kemudian $A_k$ dipetakan melalui identitas ke suku langsung $A_k$ di ruas kanan
	yang bersesuaian dengan $t'$.

	Perhatikan bahwa $t':[n]\twoheadrightarrow[k]$ berasal dari suatu pasangan
	$(i,t)$ seperti ini jika dan hanya jika $k<n$. Dengan menggunakan $\eta$
	dari Lema~\ref{prop:DK-N-unit}, definisikan $v_n$ sebagai komposisi
	$(\Gamma A)_n\xrightarrow{\text{proyeksi}}A_n
	\xrightarrow[\sim]{\eta}(\mathrm{N}\Gamma A)_n$. Uraian di atas
	menunjukkan bahwa
	% segment-id: o014.aljabr2.chapter8.dold-kan-b.q004
	\[ [\text{hasil bagi}: (\Gamma A)_n \to \Coker \Psi_n] \simeq [v_n: (\Gamma A)_n \twoheadrightarrow (\mathrm{N}\Gamma A)_n ] . \]

	Dengan menguraikan definisi, kesamaan
	$v_nu_n=\identity_{(\mathrm{N}\Gamma A)_n}$ langsung diperoleh.
	Selanjutnya akan ditunjukkan bahwa $(v_n)_n$ memberikan morfisme $v$ dalam
	$\cate{Ch}_{\geq0}(\mathcal{A})$. Ini sama artinya dengan mengatakan bahwa
	bingkai luar diagram berikut komutatif:
	% segment-id: o014.aljabr2.chapter8.dold-kan-b.g003
	\[\begin{tikzcd}[column sep=large]
		(\Gamma A)_n \arrow[r, "{\sum_i (-1)^i d^{\Gamma A}_i}"] \arrow[d, "\text{proyeksi}"'] & (\Gamma A)_{n-1} \arrow[d, "\text{proyeksi}"] \\
		A_n \arrow[r, "{\partial^A_n}"] \arrow[d, "\eta"'] & A_{n-1} \arrow[d, "\eta"] \\
		(\mathrm{N}\Gamma A)_n \arrow[r, "{\partial^{\mathrm{N}\Gamma A}_n}"'] & (\mathrm{N}\Gamma A)_{n-1}.
	\end{tikzcd}\]
	Bagian bawah komutatif karena $\eta$ adalah morfisme, sedangkan
	komutativitas bagian atas diperoleh dengan langsung menggunakan definisi
	$d_i^{\Gamma A}$ (ambil $f=\mathrm{d}^i$ dalam
	Definisi~\ref{def:DK-Gamma}). Terakhir, funktorialitas $v$ dalam $X$ jelas.
\end{proof}

% segment-id: o014.aljabr2.chapter8.dold-kan-b.p003
Tujuan berikutnya ialah membandingkan $\mathrm{C}$ dan $\mathrm{N}$ secara
lebih tajam dalam kasus kategori abelian. Morfisme antara kompleks rantai
mempunyai konsep homotopi; lihat Definisi~\ref{def:homotopy} dan
Catatan~\ref{rem:Hom-chain-cplx}. Morfisme antara kompleks rantai disebut
\emph{kuasi-isomorfisme} jika menginduksi isomorfisme pada tingkat homologi.

% segment-id: o014.aljabr2.chapter8.dold-kan-b.t002
\begin{theorem}\label{prop:Dold-Kan-N-C}
	Misalkan $\mathcal{A}$ kategori abelian dan
	$X\in\Obj(\cate{s}\mathcal{A})$. Maka
	$u:\mathrm{N}X\to\mathrm{C}X$ dari Lema~\ref{prop:Dold-Kan-u} dan
	$v:\mathrm{C}X\to\mathrm{N}X$ dari Proposisi~\ref{prop:Dold-Kan-v}
	keduanya merupakan kuasi-isomorfisme.
\end{theorem}
% segment-id: o014.aljabr2.chapter8.dold-kan-b.pf003
\begin{proof}
	Telah diketahui bahwa $vu=\identity_{\mathrm{N}X}$, sehingga cukup
	dibuktikan bahwa $v$ kuasi-isomorfisme. Karena $v$ mempunyai invers kanan,
	$v$ merupakan epimorfisme. Barisan eksak panjang bagi kompleks rantai
	(Proposisi~\ref{prop:long-exact-sequence-ses}) mereduksi masalah pada
	pembuktian bahwa $\Ker(v)$ asiklik. Berdasarkan
	Teorema~\ref{prop:Dold-Kan}, selanjutnya kita dapat mengandaikan
	$X=\Gamma A$, dengan
	$A\in\Obj(\cate{Ch}_{\geq0}(\mathcal{A}))$. Karena itu,
	% segment-id: o014.aljabr2.chapter8.dold-kan-b.q005
	\begin{gather*}
		(\mathrm{C}X)_n = \bigoplus_{t: [n] \twoheadrightarrow [k]} A_k, \\
		\partial_n : (\mathrm{C}X)_n \to (\mathrm{C}X)_{n-1}.
	\end{gather*}

	Untuk $n\geq0$ dan $i\in\Z$, tuliskan
	$(\mathrm{C}^{\leq i}X)_n$ untuk subobjek dari $(\mathrm{C}X)_n$ yang
	merupakan jumlah dari suku-suku langsung yang memenuhi
	% segment-id: o014.aljabr2.chapter8.dold-kan-b.q006
	\[ \exists \; j, \quad \max\{n-i, 1\} \leq j \leq n, \quad t(j) = t(j-1). \]
	Kita mempunyai
	$i\leq i'\implies(\mathrm{C}^{\leq i}X)_n\subset
	(\mathrm{C}^{\leq i'}X)_n$.

	Karena $v_n$ tidak lain adalah proyeksi ke suku langsung
	$t=\identity_{[n]}$, jika $i\geq n-1$ maka
	$(\mathrm{C}^{\leq i}X)_n=\Ker(v_n)$. Selanjutnya akan ditunjukkan bahwa
	% segment-id: o014.aljabr2.chapter8.dold-kan-b.e002
	\begin{equation}\label{eqn:Dold-Kan-N-C-aux0}
		\left( (\mathrm{C}^{\leq i} X)_n \right)_{n \geq 0} \;\text{membentuk subkompleks rantai dari}\; \mathrm{C}X \;\text{yang dinotasikan}\; \mathrm{C}^{\leq i} X.
	\end{equation}

	Untuk membuktikannya, misalkan diberikan pemetaan
	$t:[n]\twoheadrightarrow[k]$ dan indeks
	$\max\{n-i,1\}\leq j\leq n$ yang memenuhi kondisi di atas. Selanjutnya
	tuliskan $d_i=d_i^X$ dan $s_j=s_j^X$. Untuk setiap $0\leq h\leq n$, ambil
	$m=n-1$ dan $f=\mathrm{d}^h:[n-1]\hookrightarrow[n]$ dalam diagram
	\eqref{eqn:Dold-Kan-epi-mono}.
	\begin{itemize}
		% segment-id: o014.aljabr2.chapter8.dold-kan-b.i003
		\item Misalkan $h\notin\{j-1,j\}$ dan tetapkan
		$j':=(\mathrm{d}^h)^{-1}(j)$. Maka
		$u:[n-1]\twoheadrightarrow[l]$ memenuhi $u(j')=u(j'-1)$, dan mudah
		dilihat bahwa $n-1-i\leq j'\leq n-1$. Dalam keadaan ini,
		$(-1)^hd_h:X_n\to X_{n-1}$ memetakan $A_k$ ke suku langsung yang
		di dalam $(\mathrm{C}^{\leq i}X)_{n-1}$ ditentukan oleh $u$.
		% segment-id: o014.aljabr2.chapter8.dold-kan-b.i004
		\item Dalam keadaan pada butir sebelumnya, jika juga
		$h\geq n-i-1$, rentang $j'$ dapat dipertajam menjadi
		$n-i\leq j'\leq n-1$. Dengan kata lain, citra $A_k$ dalam hal ini
		terletak di $(\mathrm{C}^{\leq i-1}X)_{n-1}$.
		% segment-id: o014.aljabr2.chapter8.dold-kan-b.i005
		\item Untuk $h\in\{j-1,j\}$, periksa bahwa
		$t\mathrm{d}^{j-1}=t\mathrm{d}^j$. Akibatnya,
		$d_{j-1},d_j:X_n\rightrightarrows X_{n-1}$ sama pada suku langsung
		$A_k$ yang bersesuaian dengan $t$, sehingga
		$(-1)^{j-1}d_{j-1}$ dan $(-1)^jd_j$ saling meniadakan pada suku
		tersebut.
	\end{itemize}
	Hal ini membuktikan \eqref{eqn:Dold-Kan-N-C-aux0}. Butir kedua dan ketiga
	sekaligus memberikan kongruensi
	% segment-id: o014.aljabr2.chapter8.dold-kan-b.e003
	\begin{equation}\label{eqn:Dold-Kan-N-C-aux}
		\sum_{h=0}^n (-1)^h d_h \equiv \sum_{h=0}^{n-i-2} (-1)^h d_h \pmod{ \mathrm{C}^{\leq i-1} X }.
	\end{equation}

	Dengan demikian, cukup dibuktikan bahwa
	$\mathrm{C}^{\leq i}X$ asiklik. Tetapkan
	% segment-id: o014.aljabr2.chapter8.dold-kan-b.q007
	\[ D:=\mathrm{C}^{\leq i}X/\mathrm{C}^{\leq i-1}X. \]
	Perhatikan bahwa $\mathrm{C}^{<-1}X=0$. Barisan eksak panjang bagi kompleks
	rantai kemudian mereduksi masalah menjadi pembuktian secara induktif bahwa
	$D$ asiklik, untuk $i\geq0$.

	Untuk membuktikan bahwa $D$ asiklik, tinjau keluarga morfisme
	% segment-id: o014.aljabr2.chapter8.dold-kan-b.q008
	\begin{gather*}
		(-1)^{n-i-1}s_{n-i-1}:X_n\to X_{n+1}, \quad n\geq i+1.
	\end{gather*}
	Pemeriksaan langsung menunjukkan bahwa morfisme-morfisme ini mempertahankan
	$(\mathrm{C}^{\leq i-1}X)_\bullet$ dan
	$(\mathrm{C}^{\leq i}X)_\bullet$ (batas ini optimal), sehingga
	menginduksi
	% segment-id: o014.aljabr2.chapter8.dold-kan-b.q009
	\[ h_n:D_n\to D_{n+1}. \]

	Perluas definisi $s_k$ dengan nilai nol untuk setiap $k\in\Z$, lalu dengan
	cara ini perluas definisi $h_n$ dengan nilai nol ke kasus $n\leq i$.
	Perhatikan bahwa $n\leq i\implies D_n=0$, sehingga pilihan ini satu-satunya
	yang wajar. Akan ditunjukkan bahwa
	% segment-id: o014.aljabr2.chapter8.dold-kan-b.q010
	\[ \partial^D_{n+1}h_n+h_{n-1}\partial^D_n=\identity_{D_n}, \quad n\in\Z. \]
	Hal ini akan menunjukkan bahwa $D$ asiklik. Pertama, penerapan
	\eqref{eqn:Dold-Kan-N-C-aux} memberikan
	% segment-id: o014.aljabr2.chapter8.dold-kan-b.q011
	\[ \partial^D_n = \sum_{j=0}^{n-i-2} (-1)^j d_j \;\bmod\; \mathrm{C}^{\leq i-1}(X)_{n-1}. \]
	Bersama dengan \eqref{eqn:simplicial-identity}, untuk $n\geq i+1$ hal ini
	memberikan
	% segment-id: o014.aljabr2.chapter8.dold-kan-b.q012
	\begin{align*}
		\partial^D_{n+1} h_n & = (-1)^{n-i-1} \sum_{j=0}^{n-i-1} (-1)^j d_j s_{n-i-1} \; \bmod\; \mathrm{C}^{\leq i-1}(X)_n \\
		& = (-1)^{n-i-1} \sum_{j=0}^{n-i-2} (-1)^j s_{n-i-2} d_j + \identity_{\mathrm{C}^{\leq i}(X)_n} \; \bmod\; \mathrm{C}^{\leq i-1}(X)_n \\
		& = -h_{n-1} \partial^D_n + \identity_{D_n}.
	\end{align*}
	Kesamaan tersebut juga berlaku langsung untuk $n\leq i$. Pernyataan
	yang diklaim pun terbukti.
\end{proof}

% segment-id: o014.aljabr2.chapter8.dold-kan-b.r001
\begin{remark}
	Secara alami, dapat ditanyakan versi yang bersesuaian dari berbagai
	teorema dalam bagian ini, khususnya korespondensi Dold--Kan, untuk kategori
	kompleks $\cate{C}(\mathcal{A})$ atau subkategorinya
	$\cate{C}_{\leq0}(\mathcal{A})$ dan $\cate{C}_{\geq0}(\mathcal{A})$.
	Ada setidaknya dua pendekatan sederhana: pencerminan, serta pembalikan
	atau dualitas.
	\begin{enumerate}
		% segment-id: o014.aljabr2.chapter8.dold-kan-b.i006
		\item Sebagaimana dijelaskan dalam
		Catatan~\ref{rem:cochain-vs-chain}, $\cate{C}(\mathcal{A})$ dan
		$\cate{Ch}(\mathcal{A})$ saling bersesuaian melalui $X^n=X_{-n}$ dan
		$d^n=\partial_{-n}$. Semua hasil dalam bagian ini dapat dipindahkan ke
		kompleks melalui korespondensi tersebut. Sebagai contoh, versi cerminan
		Teorema~\ref{prop:Dold-Kan} adalah adjungsi yang sekaligus merupakan ekuivalensi antara
		$\cate{C}_{\leq0}(\mathcal{A})$ dan $\cate{s}\mathcal{A}$.
		% segment-id: o014.aljabr2.chapter8.dold-kan-b.i007
		\item Pendekatan lain ialah meninjau kategori objek kosimplisial
		$\cate{cs}\mathcal{A}$. Terdapat ekuivalensi kategori
		% segment-id: o014.aljabr2.chapter8.dold-kan-b.q013
		\[ \cate{cs}\mathcal{A} \simeq \cate{s}\left(\mathcal{A}^{\opp}\right)^{\opp} \simeq \cate{Ch}_{\geq 0}(\mathcal{A}^{\opp})^{\opp} \simeq \cate{C}_{\geq 0}(\mathcal{A}), \]
		dengan langkah terakhir berasal dari pembalikan anak panah, seperti dalam
		Catatan~\ref{rem:cochain-vs-chain}.
	\end{enumerate}
	Operasi seperti kompleks ternormalisasi juga dapat diterjemahkan secara
	serupa ke kategori kompleks.
\end{remark}
